\documentclass[12pt,a4paper]{article}

% ── Packages ────────────────────────────────────────────────────────────────
\usepackage[T1]{fontenc}
\usepackage[utf8]{inputenc}
\usepackage{lmodern}
\usepackage{microtype}
\usepackage{amsmath,amssymb,amsthm}
\usepackage{mathtools}
\usepackage{enumitem}
\usepackage{geometry}
\usepackage{hyperref}
\usepackage{booktabs}
\usepackage{array}
\usepackage{xcolor}
\usepackage{listings}
\usepackage{fancyhdr}

\geometry{margin=2.5cm}
\setlength{\headheight}{13.60001pt}
\addtolength{\topmargin}{-1.60001pt}

% ── Theorem environments ─────────────────────────────────────────────────────
\newtheorem{theorem}{Theorem}[section]
\newtheorem{lemma}[theorem]{Lemma}
\newtheorem{corollary}[theorem]{Corollary}
\newtheorem{proposition}[theorem]{Proposition}
\theoremstyle{definition}
\newtheorem{definition}[theorem]{Definition}
\newtheorem{remark}[theorem]{Remark}
\newtheorem{example}[theorem]{Example}

% ── Code listings ────────────────────────────────────────────────────────────
\definecolor{codegray}{rgb}{0.96,0.96,0.96}
\definecolor{codeblue}{rgb}{0.1,0.2,0.6}
\definecolor{codegreencomment}{rgb}{0.13,0.55,0.13}

\lstset{
  backgroundcolor=\color{codegray},
  basicstyle=\ttfamily\small,
  keywordstyle=\color{codeblue}\bfseries,
  commentstyle=\color{codegreencomment}\itshape,
  breaklines=true,
  frame=single,
  framerule=0pt,
  rulecolor=\color{gray!30},
  xleftmargin=6pt, xrightmargin=6pt,
}

\lstdefinelanguage{Lean4}{
  keywords={theorem,lemma,def,axiom,by,have,obtain,exact,intro,rw,push_cast,
            ring,decide,linarith,nlinarith,rcases,calc,simp,fin_cases,revert,
            import,where,fun,if,then,else,let,in,match,with,show,apply,refine,
            constructor,left,right,norm_num,norm_cast,exact_mod_cast,
            unfold,push_neg,omega,Set,Finset,open,use,private,rintro},
  sensitive=true,
  morecomment=[l]{--},
  morecomment=[s]{/-}{-/},
  morestring=[b]",
  keywordstyle=\color{codeblue}\bfseries,
  commentstyle=\color{codegreencomment}\itshape,
}

% ── Macros ───────────────────────────────────────────────────────────────────
\newcommand{\Z}{\mathbb{Z}}
\newcommand{\Q}{\mathbb{Q}}
\newcommand{\N}{\mathbb{N}}

% ── Header / Footer ──────────────────────────────────────────────────────────
\pagestyle{fancy}
\fancyhf{}
\rhead{\small\thepage}
\lhead{\small Infinitely Many Solutions to $x^5 + y^4 = 2z^2$}
\renewcommand{\headrulewidth}{0.4pt}

% ── Title ────────────────────────────────────────────────────────────────────
\title{%
  \Large\bfseries Infinitely Many Integer Solutions to\\[6pt]
  \Huge $x^5 + y^4 = 2z^2$\\[10pt]
  \large Complete Analysis with Lean~4 Formalisation
}
\author{%
  \normalsize Verified in Lean~4 / Mathlib~4 (v4.21.0)\\[2pt]
  \normalsize Repository: \href{https://github.com/JAgbanwa/miscellaneousmathproblems}%
    {\texttt{JAgbanwa/miscellaneousmathproblems}}
}
\date{\normalsize April 2026}

% ════════════════════════════════════════════════════════════════════════════
\begin{document}

\maketitle
\tableofcontents
\newpage

% ════════════════════════════════════════════════════════════════════════════
\section{Statement and Main Result}\label{sec:intro}

\subsection{The equation}

We study the Diophantine equation
\begin{equation}\label{eq:main}
  x^5 + y^4 = 2z^2, \qquad (x,y,z)\in\Z^3. \tag{$\star$}
\end{equation}

\subsection{Main theorem}

\begin{theorem}[Main result]\label{thm:main}
  The equation $x^5 + y^4 = 2z^2$ has \textbf{infinitely many} integer solutions.
  In particular, the parametric family
  \begin{equation}\label{eq:family}
    (x, y, z) = \bigl(t^4,\; t^5,\; t^{10}\bigr), \qquad t \in \Z,
  \end{equation}
  provides infinitely many distinct integer solutions.
\end{theorem}

The proof is elementary: it requires only a \texttt{ring} computation to verify
the family~\eqref{eq:family} and a strict-monotonicity argument to establish
that distinct~$t$ yield distinct solutions.  See Section~\ref{sec:proof} for the
complete proof and Section~\ref{sec:lean} for the Lean~4 formalisation.

% ════════════════════════════════════════════════════════════════════════════
\section{Derivation of the Parametric Family}\label{sec:derivation}

\subsection{Seed solution}

A direct check shows that $(x, y, z) = (1, 1, 1)$ is a solution:
\[
  1^5 + 1^4 \;=\; 1 + 1 \;=\; 2 \;=\; 2\cdot 1^2. \quad\checkmark
\]

\subsection{Weighted homogeneity}

Assign the \emph{degree weights}
\[
  \mathrm{wt}(x) = 4, \quad \mathrm{wt}(y) = 5, \quad \mathrm{wt}(z) = 10.
\]
Every monomial in~\eqref{eq:main} then has the same weighted degree:
\[
  \mathrm{wt}(x^5) = 5 \times 4 = 20, \quad
  \mathrm{wt}(y^4) = 4 \times 5 = 20, \quad
  \mathrm{wt}(2z^2) = 2 \times 10 = 20.
\]

\begin{lemma}[Scaling invariance]\label{lem:scale}
  If $(x, y, z) \in \Z^3$ satisfies~\eqref{eq:main}, then so does
  $(t^4 x,\; t^5 y,\; t^{10} z)$ for every $t \in \Z$.
\end{lemma}

\begin{proof}
  \begin{align*}
    (t^4 x)^5 + (t^5 y)^4
    &= t^{20} x^5 + t^{20} y^4
     = t^{20}(x^5 + y^4)
     = t^{20} \cdot 2z^2
     = 2(t^{10} z)^2. \qquad\square
  \end{align*}
\end{proof}

\subsection{The family}

Applying Lemma~\ref{lem:scale} to the seed $(1, 1, 1)$ with parameter $t$ yields
the family $(t^4, t^5, t^{10})$.  The key identity is:

\begin{equation}\label{eq:ring}
  (t^4)^5 + (t^5)^4
  = t^{20} + t^{20}
  = 2t^{20}
  = 2(t^{10})^2.
\end{equation}

% ════════════════════════════════════════════════════════════════════════════
\section{Proof of the Main Theorem}\label{sec:proof}

\begin{proof}[Proof of Theorem~\ref{thm:main}]
  \textbf{(Verification.)} Equation~\eqref{eq:ring} shows that every
  $(t^4, t^5, t^{10})$ is a solution.

  \textbf{(Infinitude.)} For non-negative integers $t_1 < t_2$, the
  function $t \mapsto t^4$ is strictly increasing on $\N_0$ (since $4 > 0$
  and $t \mapsto t^k$ is strictly monotone on $\N$ for $k \geq 1$).
  Hence $t_1^4 < t_2^4$, so the triples $(t_1^4, t_1^5, t_1^{10})$ and
  $(t_2^4, t_2^5, t_2^{10})$ are distinct.  Since there are infinitely many
  non-negative integers, the family yields infinitely many distinct solutions.%
  \qed
\end{proof}

% ════════════════════════════════════════════════════════════════════════════
\section{Explicit Solutions}\label{sec:explicit}

\subsection{The parametric family}

Table~\ref{tab:family} lists small members of the family $(t^4, t^5, t^{10})$.

\begin{table}[h]
\centering
\caption{Members of the parametric family $(t^4, t^5, t^{10})$.}
\label{tab:family}
\renewcommand{\arraystretch}{1.2}
\begin{tabular}{rrrrl}
\toprule
$t$ & $x = t^4$ & $y = t^5$ & $z = t^{10}$ & Verification \\
\midrule
$0$  & $0$   & $0$    & $0$      & $0 + 0 = 0 = 2 \cdot 0^2$ \\
$1$  & $1$   & $1$    & $1$      & $1 + 1 = 2 = 2 \cdot 1$ \\
$-1$ & $1$   & $-1$   & $1$      & $1 + 1 = 2 = 2 \cdot 1$ \\
$2$  & $16$  & $32$   & $1024$   & $2^{20} + 2^{20} = 2 \cdot 2^{20}$ \\
$-2$ & $16$  & $-32$  & $1024$   & same \\
$3$  & $81$  & $243$  & $59\,049$ & $3^{20} + 3^{20} = 2 \cdot 3^{20}$ \\
$4$  & $256$ & $1024$ & $1\,048\,576$ & $4^{20} + 4^{20} = 2 \cdot 4^{20}$ \\
\bottomrule
\end{tabular}
\end{table}

\subsection{Secondary family ($y = 0$)}

Setting $y = 0$ reduces~\eqref{eq:main} to $x^5 = 2z^2$.  Writing $x = 2m^2$:
\[
  (2m^2)^5 = 2^5 m^{10} = 2 \cdot (4m^5)^2.
\]
Hence $(x,y,z) = (2m^2,\; 0,\; 4m^5)$ is a solution for every $m \in \Z$.

\begin{example}
  $m = 1$: $(x,y,z) = (2, 0, 4)$.  Check: $2^5 + 0 = 32 = 2 \cdot 16 = 2 \cdot 4^2$. $\checkmark$
\end{example}

\subsection{Third family ($z = 0$)}

Setting $z = 0$ reduces~\eqref{eq:main} to $x^5 + y^4 = 0$, i.e.\ $x^5 = -y^4$.

Since $y^4 \geq 0$ we need $x \leq 0$.  Write $x = -a$, $a \geq 0$; then $a^5 = y^4$.
By unique factorisation, writing the $p$-adic valuation condition $5v_p(a) = 4v_p(y)$
forces $4 \mid v_p(a)$ for every prime $p$, so $a = n^4$ and $y = n^5$ for some $n \in \N_0$.
This identifies the complete $z = 0$ family:

\begin{equation}\label{eq:family3}
  (x, y, z) = \bigl(-k^4,\; k^5,\; 0\bigr), \qquad k \in \Z.
\end{equation}

\begin{proof}[Verification]
  $(-k^4)^5 + (k^5)^4 = -k^{20} + k^{20} = 0 = 2\cdot 0^2.$
\end{proof}

Note that the weighted-homogeneity scaling by $t$ sends $(-1, 1, 0) \mapsto (-t^4, t^5, 0)$,
recovering the entire family from the single primitive seed $(-1, 1, 0)$.

% ════════════════════════════════════════════════════════════════════════════
\section{Primitive Integer Solutions}\label{sec:primitive}

\begin{definition}
  A solution $(x, y, z) \in \Z^3$ is \emph{primitive} if
  $\gcd(|x|,|y|,|z|) = 1$ (taking the gcd of the nonzero entries when one
  coordinate is zero).
\end{definition}

\subsection{Primitivity of each parametric family}

\begin{proposition}\label{prop:prim-families}
  \begin{enumerate}[label=\emph{(\roman*)}]
    \item \textbf{Family 1} $(t^4, t^5, t^{10})$:
      $\gcd(t^4, t^5, t^{10}) = t^4$.  Primitive if and only if $|t| = 1$.
    \item \textbf{Family 2} $(2m^2, 0, 4m^5)$:
      $\gcd(2m^2, 4|m|^5) = 2m^2 \geq 2$ for $m \neq 0$.  Never primitive.
    \item \textbf{Family 3} $(-k^4, k^5, 0)$:
      $\gcd(k^4, |k|^5) = k^4$.  Primitive if and only if $|k| = 1$.
  \end{enumerate}
\end{proposition}

\subsection{Complete list of primitive solutions}

\begin{theorem}[Primitive solutions]\label{thm:primitive}
  A brute-force search over $|x|, |y| \leq 500$ shows that
  all primitive (i.e.\ $\gcd = 1$) solutions are:

  \begin{table}[h]
  \centering
  \caption{Complete list of primitive solutions (up to $|x|, |y| \leq 500$).}
  \label{tab:primitive}
  \renewcommand{\arraystretch}{1.2}
  \begin{tabular}{rrrl}
  \toprule
  $x$ & $y$ & $z$ & Source \\
  \midrule
  $1$  & $1$  & $1$  & Family 1, $t = 1$ \\
  $1$  & $1$  & $-1$ & $z \to -z$ \\
  $1$  & $-1$ & $1$  & Family 1, $t = -1$ \\
  $1$  & $-1$ & $-1$ & $z \to -z$ \\
  $-1$ & $1$  & $0$  & Family 3, $k = 1$ \\
  $-1$ & $-1$ & $0$  & Family 3, $k = -1$ \\
  \bottomrule
  \end{tabular}
  \end{table}

  All primitive solutions arise from exactly two weighted-primitive seeds:
  \begin{align}
    \text{Seed A:} &\quad (1,\; 1,\; 1)
      \xrightarrow{\;\text{scale by }t\;}
      (t^4,\; t^5,\; t^{10}), \\[4pt]
    \text{Seed B:} &\quad (-1,\; 1,\; 0)
      \xrightarrow{\;\text{scale by }t\;}
      (-t^4,\; t^5,\; 0).
  \end{align}
\end{theorem}

\begin{remark}
  The $z \to -z$ symmetry is trivial since $2z^2 = 2(-z)^2$.
  The $y \to -y$ symmetry reflects the fact that $y^4 = (-y)^4$.
  Taking these symmetries into account, there are exactly \emph{two}
  primitive solutions up to sign changes: $(1, 1, 1)$ and $(-1, 1, 0)$.
\end{remark}

% ════════════════════════════════════════════════════════════════════════════
\section{Structural Observations}\label{sec:structure}

\subsection{Weighted projective interpretation}

The equation~\eqref{eq:main} defines a hypersurface in the
\emph{weighted projective space} $\mathbb{P}^2_{(4,5,10)}$ (weights
$4, 5, 10$, common weight $20$).  The rational point
$[1:1:1]_{(4,5,10)}$ gives rise to the entire parametric family when
pulled back to affine coordinates via $t \mapsto [t^4 : t^5 : t^{10}]$.

\subsection{Why Faltings' theorem does not apply}

Faltings' theorem (every smooth projective curve of genus $g \geq 2$
over $\Q$ has finitely many rational points) does not apply here
because the `curve' swept out by the parametric family is a rational
curve (parametrised by one free variable $t$), which has genus $0$.
Accordingly, infinitely many integral points are entirely expected
and confirmed.

\subsection{Parity constraint}

\begin{proposition}\label{prop:parity}
  Any integer solution $(x, y, z)$ to~\eqref{eq:main} satisfies $x \equiv y \pmod{2}$.
\end{proposition}

\begin{proof}
  Since $2z^2 \equiv 0 \pmod{2}$, we need $x^5 + y^4 \equiv 0 \pmod{2}$.
  Now $x^5 \equiv x \pmod{2}$ and $y^4 \equiv y^2 \equiv y \cdot y \equiv y^2 \pmod 2$;
  more precisely $y^4 \equiv 0$ if $y$ is even and $\equiv 1$ if $y$ is odd.
  Similarly $x^5 \equiv 0$ if $x$ is even and $\equiv 1$ if $x$ is odd.
  The sum is $\equiv 0 \pmod 2$ iff $x$ and $y$ have the same parity.%
  \qed
\end{proof}

The parametric family satisfies this: $t^4$ and $t^5$ always have the same
parity (both even iff $t$ is even, both odd iff $t$ is odd).

% ════════════════════════════════════════════════════════════════════════════
\section{Lean 4 Formalisation}\label{sec:lean}

The file \texttt{InfinitelyManySolutions.lean} contains a complete Lean~4 /
Mathlib~4 formalisation of Theorem~\ref{thm:main}.

\textbf{Sorry count: 0.  Axiom count: 0} (beyond Lean's standard foundations and
Mathlib).

\begin{table}[h]
\centering
\caption{Lean~4 proof components.}
\label{tab:lean}
\renewcommand{\arraystretch}{1.2}
\begin{tabular}{lll}
\toprule
\textbf{Lean name} & \textbf{Statement} & \textbf{Method} \\
\midrule
\texttt{parametric\_solution} & $(t^4)^5 + (t^5)^4 = 2(t^{10})^2$ for all $t : \Z$ & \texttt{ring} \\
\texttt{natMap\_mem} & Every $(n^4, n^5, n^{10})$ is a solution & \texttt{simp only}, \texttt{ring} \\
\texttt{pow4\_strictMono} & $n \mapsto n^4$ is strictly mono on $\N$ & \texttt{Nat.pow\_lt\_pow\_left} \\
\texttt{natMap\_injective} & $n \mapsto (n^4, n^5, n^{10})$ is injective & strict monotonicity \\
\texttt{solutions\_infinite} & The solution set is infinite & \texttt{Set.infinite\_range\_of\_injective} \\
\bottomrule
\end{tabular}
\end{table}

\subsection{Core Lean code}

\begin{lstlisting}[language=Lean4]
theorem parametric_solution (t : Z) :
    isolution (t ^ 4) (t ^ 5) (t ^ 10) := by
  unfold isolution
  ring

private theorem pow4_strictMono : StrictMono (fun n : N => n ^ 4) := by
  intro a b hab
  exact Nat.pow_lt_pow_left hab (by norm_num)

theorem natMap_injective : Function.Injective natMap := by
  intro n1 n2 h
  simp only [natMap, Prod.mk.injEq] at h
  have h4 : n1 ^ 4 = n2 ^ 4 := by exact_mod_cast h.1
  exact pow4_strictMono.injective h4

theorem natMap_mem (n : N) : natMap n \in solutions := by
  simp only [solutions, Set.mem_setOf_eq, isolution, natMap]
  ring

theorem solutions_infinite : solutions.Infinite := by
  apply (Set.infinite_range_of_injective natMap_injective).mono
  rintro _ <n, rfl>
  exact natMap_mem n
\end{lstlisting}

The proof is entirely mechanical: equation verification by term rewriting
(\texttt{ring}), injectivity by standard monotonicity lemmas, and infinitude
by the Mathlib utility \texttt{Set.infinite\_range\_of\_injective}.  No
advanced arithmetic geometry is required.

% ════════════════════════════════════════════════════════════════════════════
\section{Summary}\label{sec:summary}

\begin{theorem}[Summary]
  The Diophantine equation $x^5 + y^4 = 2z^2$ has infinitely many integer
  solutions.  The equation admits three known parametric families:
  \begin{enumerate}[label=\emph{(\roman*)}]
    \item $(t^4, t^5, t^{10})$ for $t \in \Z$ — the principal family, from the seed $(1,1,1)$.
    \item $(2m^2, 0, 4m^5)$ for $m \in \Z$ — the $y=0$ branch.
    \item $(-k^4, k^5, 0)$ for $k \in \Z$ — the $z=0$ branch, from the seed $(-1,1,0)$.
  \end{enumerate}
  The \textbf{primitive solutions} ($\gcd = 1$, up to signs of $y$ and $z$)
  are exactly $(1, 1, 1)$ and $(-1, 1, 0)$ (within the search bound $|x|,|y|\leq 500$).
\end{theorem}

\begin{table}[h]
\centering
\caption{Proof status summary.}
\label{tab:status}
\renewcommand{\arraystretch}{1.2}
\begin{tabular}{lll}
\toprule
\textbf{Component} & \textbf{Status} & \textbf{Method} \\
\midrule
Seed solution $(1,1,1)$ & \checkmark\ Complete & Direct substitution \\
Parametric family verified & \checkmark\ Complete & \texttt{ring} identity \\
$z=0$ family $(-k^4, k^5, 0)$ & \checkmark\ Complete & Unique factorisation \\
Infinitely many distinct solutions & \checkmark\ Complete & Strict monotonicity of $n^4$ \\
Primitive solutions (brute force, $\leq 500$) & \checkmark\ Complete & Two seeds: $(1,1,1)$, $(-1,1,0)$ \\
Lean~4 formalisation (0 sorry) & \checkmark\ Complete & Mathlib tactics \\
Classification of \emph{all} solutions & Open & Weighted projective geometry \\
\bottomrule
\end{tabular}
\end{table}

\end{document}
